\chapter{Lazy eye (Amblyopia)}

\section*{Definition and Overview}
Lazy eye (amblyopia) is a childhood vision development disorder where one eye fails to achieve normal visual acuity, even with glasses or contact lenses. It occurs when the brain and the affected eye do not work together properly, leading the brain to favour the other eye and actively suppress the visual input from the weaker eye. Amblyopia is the leading cause of single-eye vision loss in children and young adults, but it is highly treatable if diagnosed and addressed early in life, typically before the age of 7 or 8.

\section*{Epidemiology}
Amblyopia affects approximately 2 to 3\% of the global population, making it the most common cause of visual impairment in children. It typically develops in early infancy or childhood. There is no major sex or racial predilection. In Australia, child health checks and preschool screening programmes are designed specifically to detect amblyopia early, as late diagnosis significantly reduces treatment efficacy.

\section*{Aetiology and Pathophysiology}
Amblyopia arises when normal visual development is disrupted during the critical period of brain development in childhood:

\begin{itemize}
    \item \textbf{Strabismus:} Misalignment of the eyes (squint) causes double vision, prompting the brain to ignore the input from the turned eye.
    \item \textbf{Refractive errors:} A significant difference in prescription between the two eyes (anisometropia) or high refractive errors in both eyes cause one or both images to be chronically blurred.
    \item \textbf{Deprivation:} Obstruction of the visual axis, such as from a congenital cataract, ptosis (drooping eyelid), or corneal opacity, prevents clear light and images from reaching the retina.
    \item \textbf{Cortical suppression:} The visual cortex of the brain suppresses the blurred or misaligned input from the affected eye, resulting in permanent developmental deficit of the corresponding neural pathways if left untreated.
\end{itemize}

\section*{Clinical Features}
Amblyopia is often asymptomatic and can go unnoticed without screening, but some signs may be present in a child:

\begin{itemize}
    \item \textbf{Poor depth perception:} Difficulty judging distances, leading to clumsiness, frequent bumping into objects, or trouble catching balls.
    \item \textbf{Squinting or head tilting:} The child may tilt their head, turn their face, or squint to see more clearly or compensate for misalignment.
    \item \textbf{Eye misalignment:} A visible squint (strabismus) where one eye turns inwards, outwards, upwards, or downwards.
    \item \textbf{Difficulty with close tasks:} Trouble with reading, writing, or colouring, or covering one eye when trying to focus.
\end{itemize}

\section*{Differential Diagnosis}
Other conditions that reduce vision in children must be ruled out:

\begin{itemize}
    \item \textbf{Organic vision loss:} Structural eye diseases such as optic nerve hypoplasia, retinal dystrophy, or intraocular tumours.
    \item \textbf{Uncorrected refractive error:} Simple nearsightedness, farsightedness, or astigmatism that improves immediately and fully with glasses.
    \item \textbf{Cortical visual impairment:} Vision loss due to brain injury rather than developmental suppression of the eye.
    \item \textbf{Congenital cataracts:} Opacity of the lens that must be surgically cleared.
\end{itemize}

\section*{When to Seek Medical Care}
Parents should arrange an eye examination for their child if they notice an eye turn, frequent squinting, head tilting, or if the child complains of difficulty seeing. Universal vision screening is recommended for all children between 3 and 5 years of age.

\textbf{Call 000 (or local emergency services) for:}
\begin{itemize}
    \item Sudden loss of vision in one or both eyes
    \item Severe eye pain, swelling, or redness
    \item A white reflex (leukocoria) in the child's pupil, which could indicate a serious condition like retinoblastoma
    \item Sudden onset of double vision or a new squint following a head injury
    \item Severe drowsiness, confusion, or vomiting accompanying vision changes
\end{itemize}

\section*{Diagnosis and Investigations}
A comprehensive paediatric eye assessment is used to diagnose amblyopia:

\begin{itemize}
    \item \textbf{Visual acuity testing:} Age-appropriate charts (e.g., symbols, shapes, or letters) to measure and compare vision in each eye.
    \item \textbf{Refraction testing:} Using eye drops (cycloplegic drops) to dilate the pupils and temporarily relax focusing muscles, allowing accurate measurement of refractive errors.
    \item \textbf{Cover test:} Covering one eye at a time to check for alignment and detect strabismus.
    \item \textbf{Ophthalmoscopy:} Examining the back of the eye to ensure the retina and optic nerve are structurally healthy.
\end{itemize}

\section*{Management}
Treatment aims to correct underlying causes and force the brain to use the weaker eye:

\begin{itemize}
    \item \textbf{Corrective eyewear:} Prescribing glasses to correct refractive differences and provide clear images to both retinas.
    \item \textbf{Occlusion therapy (patching):} Placing an opaque patch over the stronger eye for several hours a day to stimulate visual development in the lazy eye.
    \item \textbf{Atropine eye drops:} Instilling drops in the stronger eye to temporarily blur near vision, forcing the child to use the weaker eye.
    \item \textbf{Surgery:} Correcting strabismus (squint surgery) or removing cataracts to clear the visual path, usually followed by patching.
    \item \textbf{Vision therapy:} Specialized exercises to improve eye coordination, tracking, and binocular vision.
\end{itemize}

\section*{Complications}
\begin{itemize}
    \item Permanent, irreversible loss of visual acuity in the affected eye if not treated before the visual pathways mature (around age 8)
    \item Loss of binocular vision and stereopsis (depth perception), affecting future career options or sports performance
    \item Higher vulnerability to total blindness if the healthy eye is injured or diseased later in life
    \item Poor school performance or developmental delays due to undetected visual impairment
    \item Psychological issues, such as low self-esteem or bullying associated with wearing patches or glasses
\end{itemize}

\section*{Prognosis and Outcomes}
The prognosis is excellent if treatment is started early in childhood. Most children achieve normal or near-normal vision in the lazy eye if compliant with patching or atropine drops. However, treatment efficacy drops significantly after age 7, and is rarely successful after age 10 to 12. Early detection via screening is the key predictor of a successful outcome.

\section*{Prevention and Self-Care}
\begin{itemize}
    \item Ensure children undergo routine vision screening by a GP, paediatrician, or maternal health nurse between ages 3 and 5
    \item Follow the prescribed patching schedule strictly and make it a fun, rewarding daily routine for the child
    \item Monitor the child to ensure they are not peeking around the edges of the eye patch
    \item Ensure glasses are worn consistently and adjusted regularly for a comfortable fit
    \item Encourage activities that require hand-eye coordination, such as colouring or puzzles, while the patch is on
\end{itemize}

\section*{Sources and Further Reading}
The following organisations provide reliable information on amblyopia:

\begin{itemize}
    \item RANZCO. \textit{Patient information on amblyopia and paediatric eye health.} https://ranzco.edu/
    \item NHS. \textit{Overview of lazy eye causes, signs, and treatments.} https://www.nhs.uk/
    \item Mayo Clinic. \textit{Clinical guide to amblyopia diagnosis and patching.} https://www.mayoclinic.org/
    \item Healthdirect Australia. \textit{Lazy eye (amblyopia) symptoms and support.} https://www.healthdirect.gov.au/
    \item American Association for Pediatric Ophthalmology and Strabismus. \textit{Educational resources for parents.} https://aapos.org/
\end{itemize}
